\documentclass[11pt]{article}

% ============================================================
% Packages
% ============================================================
\usepackage[a4paper,margin=1in]{geometry}
\usepackage{amsmath,amssymb,amsthm,mathtools}
\usepackage{bm}
\usepackage{enumitem}
\usepackage{microtype}
\usepackage{hyperref}

% ============================================================
% Theorem environments
% ============================================================
\newtheorem{theorem}{Theorem}[section]
\newtheorem{proposition}[theorem]{Proposition}
\newtheorem{lemma}[theorem]{Lemma}
\newtheorem{corollary}[theorem]{Corollary}
\newtheorem{claim}[theorem]{Claim}
\newtheorem{definition}[theorem]{Definition}
\newtheorem{remark}[theorem]{Remark}

% ============================================================
% Commands
% ============================================================
\newcommand{\R}{\mathbb{R}}
\newcommand{\F}{\mathrm{F}}
\newcommand{\op}{\mathrm{op}}
\newcommand{\ip}[2]{\left\langle #1,#2\right\rangle}
\newcommand{\norm}[1]{\left\lVert #1\right\rVert}
\newcommand{\abs}[1]{\left\lvert #1\right\rvert}
\newcommand{\Sph}{\mathbb{S}}
\newcommand{\cS}{\mathcal{S}}
\newcommand{\cN}{\mathcal{N}}
\newcommand{\cC}{\mathcal{C}}
\newcommand{\E}{\mathbb{E}}
\newcommand{\Pp}{\mathbb{P}}
\newcommand{\1}{\mathbf{1}}
\newcommand{\argmin}{\operatorname*{arg\,min}}
\newcommand{\Var}{\operatorname{Var}}
\newcommand{\vecop}{\operatorname{vec}}
\newcommand{\Span}{\operatorname{span}}
\newcommand{\diag}{\operatorname{diag}}

% ============================================================
% Title
% ============================================================
\title{
A Technical Note on Existence and Uniqueness\\
of the Finite-Population Softmax Associative-Memory Minimizer
}
\author{}
\date{}

\begin{document}

\maketitle

\begin{abstract}
We study the finite-population softmax loss
\[
L(W)
=
\sum_{i=1}^{N}p_i
\left(
\log\sum_{j=1}^{N}
\exp(u_j^\top Wv_i)
-
u_i^\top Wv_i
\right),
\qquad
W\in\R^{d\times d},
\]
where all weights satisfy \(p_i>0\), and
\[
u_i,v_i\overset{\mathrm{iid}}{\sim}N(0,I_d/d),
\]
with the two Gaussian families mutually independent.

The objective need not possess a finite minimizer for every realization:
if there exists a nonzero matrix direction along which every correct class
weakly dominates every incorrect class, the loss decreases toward its
infimum. We prove that when \(N>d^2\), such a
direction fails to exist with probability at least
\[
1-
2^{1-N}
\sum_{k=0}^{d^2-1}
\binom{N-1}{k}.
\]
On the same event the loss is coercive. Gaussian spanning implies strict
convexity almost surely, and consequently the minimizer exists, is finite,
and is unique.

For \(N=\lceil d^\beta\rceil\) with fixed \(\beta>2\), the exceptional
probability is \(\exp[-\Omega(N)]\).
\end{abstract}

\tableofcontents

% ============================================================
\section{Model and main theorem}
% ============================================================

Fix integers
\[
d\ge 2,
\qquad
N>d^2.
\]
Let
\[
u_1,\dots,u_N,
\qquad
v_1,\dots,v_N
\]
be random vectors in \(\R^d\) satisfying
\[
u_i,v_i\overset{\mathrm{iid}}{\sim}N(0,I_d/d),
\]
where the family \((u_i)_{i=1}^N\) is independent of the family
\((v_i)_{i=1}^N\).

Let
\[
p_1,\dots,p_N>0,
\qquad
\sum_{i=1}^N p_i=1.
\]
The proof below only uses strict positivity of the weights. In particular,
it applies to the power-law distribution
\[
p_i
=
\frac{i^{-\alpha}}
{\sum_{k=1}^N k^{-\alpha}},
\qquad
\alpha\ge 0.
\]

For \(W\in\R^{d\times d}\), define the logits
\[
s_{ji}(W)=u_j^\top Wv_i,
\qquad
1\le i,j\le N,
\]
the per-item loss
\[
\ell_i(W)
=
\log\left(
\sum_{j=1}^{N}e^{s_{ji}(W)}
\right)
-
s_{ii}(W),
\]
and the total loss
\[
L(W)
=
\sum_{i=1}^{N}p_i\ell_i(W).
\]

Equivalently,
\begin{equation}
\ell_i(W)
=
\log\left(
\sum_{j=1}^{N}
\exp\bigl((u_j-u_i)^\top Wv_i\bigr)
\right).
\label{eq:relative-loss}
\end{equation}

We equip \(\R^{d\times d}\) with the Frobenius inner product
\[
\ip{A}{B}
=
\operatorname{Tr}(A^\top B)
\]
and Frobenius norm
\[
\norm{A}_{\F}
=
\sqrt{\ip{A}{A}}.
\]

\begin{theorem}[Finite existence and uniqueness]
\label{thm:main}
Assume \(N>d^2\). Then
\begin{equation}
\Pp\left(
L\text{ has a unique finite minimizer }W_\star
\right)
\ge
1-
2^{1-N}
\sum_{k=0}^{d^2-1}
\binom{N-1}{k}.
\label{eq:main-probability}
\end{equation}

Equivalently, except on an event of probability at most
\[
2^{1-N}
\sum_{k=0}^{d^2-1}
\binom{N-1}{k},
\]
there exists a unique matrix \(W_\star\in\R^{d\times d}\) such that
\[
L(W_\star)
=
\inf_{W\in\R^{d\times d}}L(W).
\]

In particular, if
\[
N=\left\lceil d^\beta\right\rceil
\]
for a fixed \(\beta>2\), then, for all sufficiently large \(d\),
\begin{equation}
\Pp\left(
L\text{ has a unique finite minimizer}
\right)
\ge
1-\exp\left(-\frac{N-1}{16}\right).
\label{eq:main-exponential}
\end{equation}
Consequently,
\[
\Pp\left(
L\text{ has a unique finite minimizer}
\right)
=
1-\exp[-\Omega(N)].
\]
\end{theorem}

The remainder of this note proves Theorem~\ref{thm:main} from first
principles.

% ============================================================
\section{Deterministic geometry of the loss}
% ============================================================

\subsection{The recession functional}

For \(H\in\R^{d\times d}\), define
\begin{equation}
L^\infty(H)
=
\sum_{i=1}^{N}p_i
\left[
\max_{1\le j\le N}u_j^\top Hv_i
-
u_i^\top Hv_i
\right].
\label{eq:recession}
\end{equation}

The notation \(L^\infty\) is justified because this is the recession
function of \(L\). For the existence proof, we only need its explicit
formula, continuity, and positive homogeneity.

\begin{lemma}[Basic properties of the recession functional]
\label{lem:recession-basic}
For every \(H\in\R^{d\times d}\):

\begin{enumerate}[label=(\roman*)]
\item \(L^\infty(H)\ge 0\);

\item for every \(t\ge 0\),
\[
L^\infty(tH)=tL^\infty(H);
\]

\item \(L^\infty\) is continuous;

\item for every \(W\in\R^{d\times d}\),
\begin{equation}
L(W)\ge L^\infty(W).
\label{eq:L-lower-recession}
\end{equation}
\end{enumerate}
\end{lemma}

\begin{proof}
For each fixed \(i\), the index \(j=i\) is included in the maximum.
Therefore
\[
\max_j u_j^\top Hv_i-u_i^\top Hv_i
\ge
u_i^\top Hv_i-u_i^\top Hv_i
=
0.
\]
Since \(p_i>0\), summing proves \(L^\infty(H)\ge 0\).

For \(t\ge 0\),
\[
\max_j u_j^\top(tH)v_i-u_i^\top(tH)v_i
=
t
\left(
\max_j u_j^\top Hv_i-u_i^\top Hv_i
\right),
\]
which proves positive homogeneity.

Each map
\[
H\longmapsto u_j^\top Hv_i
\]
is linear and therefore continuous. A maximum of finitely many continuous
functions is continuous. Hence each summand in
\eqref{eq:recession} is continuous, and so is \(L^\infty\).

Finally, for every collection of real numbers \(a_1,\dots,a_N\),
\[
\log\sum_{j=1}^{N}e^{a_j}
\ge
\max_j a_j.
\]
Applying this with
\[
a_j=u_j^\top Wv_i
\]
gives
\[
\ell_i(W)
=
\log\sum_j e^{u_j^\top Wv_i}
-
u_i^\top Wv_i
\ge
\max_j u_j^\top Wv_i-u_i^\top Wv_i.
\]
Multiplying by \(p_i\) and summing over \(i\) yields
\eqref{eq:L-lower-recession}.
\end{proof}

\subsection{The weak-separator cone}

Define the weak-separator cone
\begin{equation}
\cS
=
\left\{
H\in\R^{d\times d}:
u_i^\top Hv_i
\ge
u_j^\top Hv_i
\text{ for every }i,j
\right\}.
\label{eq:separator-cone}
\end{equation}

Thus \(H\in\cS\) means that, in the direction \(H\), the correct candidate
\(u_i\) weakly dominates every competing candidate \(u_j\) for every
query \(v_i\).

Also define the lineality space
\begin{equation}
\cN
=
\left\{
H\in\R^{d\times d}:
(u_j-u_1)^\top Hv_i=0
\text{ for every }i,j
\right\}.
\label{eq:lineality}
\end{equation}

A direction \(H\in\cN\) adds the same logit to all candidates in each
column and therefore leaves the loss unchanged.

\begin{lemma}[Zero set of \(L^\infty\)]
\label{lem:zero-set}
For every \(H\in\R^{d\times d}\),
\[
L^\infty(H)=0
\quad\Longleftrightarrow\quad
H\in\cS.
\]
\end{lemma}

\begin{proof}
Every summand in \eqref{eq:recession} is nonnegative. Since every
\(p_i>0\), the sum vanishes if and only if each summand vanishes:
\[
\max_j u_j^\top Hv_i-u_i^\top Hv_i=0
\qquad
\text{for every }i.
\]
This is equivalent to
\[
u_i^\top Hv_i
\ge
u_j^\top Hv_i
\qquad
\text{for every }i,j,
\]
which is exactly \(H\in\cS\).
\end{proof}

\subsection{Why a nonzero separator prevents finite minimization}

\begin{lemma}[Strict descent along a nontrivial separator]
\label{lem:strict-descent}
Assume \(\cN=\{0\}\). If
\[
H\in\cS\setminus\{0\},
\]
then, for every fixed \(W\in\R^{d\times d}\), the function
\[
t\longmapsto L(W+tH)
\]
is strictly decreasing on \(\R\).

Consequently, \(L\) cannot possess a finite minimizer.
\end{lemma}

\begin{proof}
Fix \(W\) and \(H\in\cS\). Define
\[
\delta_{ji}
=
u_i^\top Hv_i-u_j^\top Hv_i.
\]
Because \(H\in\cS\),
\[
\delta_{ji}\ge 0
\qquad
\text{for every }i,j,
\]
and
\[
\delta_{ii}=0.
\]

Using the relative-logit representation
\eqref{eq:relative-loss},
\begin{align}
\ell_i(W+tH)
&=
\log
\sum_{j=1}^{N}
\exp\left(
(u_j-u_i)^\top(W+tH)v_i
\right)
\nonumber\\
&=
\log
\sum_{j=1}^{N}
\exp\left(
(u_j-u_i)^\top Wv_i
-
t\delta_{ji}
\right).
\label{eq:loss-along-separator}
\end{align}

Every term inside the sum in \eqref{eq:loss-along-separator} is
nonincreasing in \(t\). Therefore every
\[
t\longmapsto \ell_i(W+tH)
\]
is nonincreasing.

We now show that at least one per-item loss is strictly decreasing.
Suppose, to the contrary, that
\[
\delta_{ji}=0
\qquad
\text{for every }i,j.
\]
Then
\[
u_i^\top Hv_i=u_j^\top Hv_i
\qquad
\text{for every }i,j.
\]
In particular,
\[
(u_j-u_1)^\top Hv_i=0
\qquad
\text{for every }i,j,
\]
so \(H\in\cN\). Since \(\cN=\{0\}\), this would imply \(H=0\), contrary
to the assumption.

Hence there exist \(i_0,j_0\) such that
\[
\delta_{j_0i_0}>0.
\]
For this fixed \(i_0\), at least one positive term in the exponential
sum is strictly decreasing in \(t\), while no term increases. Thus
\[
t\longmapsto \ell_{i_0}(W+tH)
\]
is strictly decreasing.

Because \(p_{i_0}>0\), the weighted sum
\[
L(W+tH)=\sum_i p_i\ell_i(W+tH)
\]
is strictly decreasing.

If \(W_\star\) were a finite minimizer, then for every \(t>0\),
\[
L(W_\star+tH)<L(W_\star),
\]
contradicting minimality. Hence no finite minimizer exists.
\end{proof}

\subsection{Coercivity when no separator exists}

\begin{proposition}[Deterministic existence criterion]
\label{prop:deterministic-existence}
Assume
\[
\cS=\{0\}.
\]
Then \(L\) is coercive in Frobenius norm: there exists
\(\kappa>0\) such that
\begin{equation}
L(W)\ge \kappa\norm{W}_{\F}
\qquad
\text{for every }W\in\R^{d\times d}.
\label{eq:coercive-bound}
\end{equation}
Consequently, \(L\) attains its minimum at some finite matrix
\(W_\star\in\R^{d\times d}\).
\end{proposition}

\begin{proof}
Consider the Frobenius unit sphere
\[
\Sph_{\F}
=
\left\{
H\in\R^{d\times d}:\norm{H}_{\F}=1
\right\}.
\]
This is compact because \(\R^{d\times d}\) is finite-dimensional.

By Lemma~\ref{lem:recession-basic}, \(L^\infty\) is continuous.
By Lemma~\ref{lem:zero-set} and the assumption \(\cS=\{0\}\),
\[
L^\infty(H)>0
\qquad
\text{for every }H\in\Sph_{\F}.
\]
A continuous, strictly positive function on a compact set attains a
strictly positive minimum. Define
\[
\kappa
=
\min_{H\in\Sph_{\F}}L^\infty(H)>0.
\]

Let \(W\neq 0\). By positive homogeneity,
\[
L^\infty(W)
=
\norm{W}_{\F}
L^\infty\left(
\frac{W}{\norm{W}_{\F}}
\right)
\ge
\kappa\norm{W}_{\F}.
\]
For \(W=0\), the same inequality holds trivially. Using
\eqref{eq:L-lower-recession},
\[
L(W)
\ge
L^\infty(W)
\ge
\kappa\norm{W}_{\F},
\]
which proves \eqref{eq:coercive-bound}.

It remains to prove attainment. Let
\[
m=\inf_{W\in\R^{d\times d}}L(W).
\]
Since \(L(0)=\log N<\infty\), we have \(m<\infty\).
Choose a minimizing sequence \((W_k)\) satisfying
\[
L(W_k)\longrightarrow m.
\]
For all sufficiently large \(k\),
\[
L(W_k)\le m+1.
\]
By \eqref{eq:coercive-bound},
\[
\kappa\norm{W_k}_{\F}
\le
L(W_k)
\le
m+1.
\]
Hence \((W_k)\) is bounded in the finite-dimensional space
\(\R^{d\times d}\). It therefore has a convergent subsequence,
still denoted \(W_k\), such that
\[
W_k\longrightarrow W_\star.
\]
The function \(L\) is continuous, because it is a finite sum of
compositions of continuous functions. Consequently,
\[
L(W_\star)
=
\lim_{k\to\infty}L(W_k)
=
m.
\]
Thus \(W_\star\) is a finite minimizer.
\end{proof}

% ============================================================
\section{Strict convexity and uniqueness}
% ============================================================

For \(W\in\R^{d\times d}\), define the softmax probabilities
\begin{equation}
\pi_{j\mid i}(W)
=
\frac{
\exp(u_j^\top Wv_i)
}{
\sum_{k=1}^{N}\exp(u_k^\top Wv_i)
}.
\label{eq:softmax}
\end{equation}
For every finite \(W\),
\[
\pi_{j\mid i}(W)>0
\qquad
\text{for every }i,j.
\]

\begin{lemma}[Gradient and Hessian in a matrix direction]
\label{lem:hessian}
For every \(W,H\in\R^{d\times d}\),
\begin{align}
\frac{d}{dt}L(W+tH)\bigg|_{t=0}
&=
\sum_{i=1}^{N}p_i
\left[
\sum_{j=1}^{N}
\pi_{j\mid i}(W)\,
u_j^\top Hv_i
-
u_i^\top Hv_i
\right],
\label{eq:first-directional}
\\
\frac{d^2}{dt^2}L(W+tH)\bigg|_{t=0}
&=
\sum_{i=1}^{N}p_i
\Var_{j\sim\pi_{\cdot\mid i}(W)}
\left(
u_j^\top Hv_i
\right).
\label{eq:second-directional}
\end{align}
\end{lemma}

\begin{proof}
Fix \(i\), and define
\[
a_{ji}=u_j^\top Hv_i,
\qquad
b_{ji}=u_j^\top Wv_i.
\]
Then
\[
\ell_i(W+tH)
=
\log\sum_{j=1}^{N}e^{b_{ji}+ta_{ji}}
-
b_{ii}
-
ta_{ii}.
\]

Differentiating,
\begin{align*}
\frac{d}{dt}\ell_i(W+tH)
&=
\frac{
\sum_j a_{ji}e^{b_{ji}+ta_{ji}}
}{
\sum_k e^{b_{ki}+ta_{ki}}
}
-
a_{ii}.
\end{align*}
At \(t=0\), this becomes
\[
\sum_j\pi_{j\mid i}(W)a_{ji}-a_{ii}.
\]
Multiplying by \(p_i\) and summing proves
\eqref{eq:first-directional}.

For the second derivative, let
\[
\pi_{j\mid i}(t)
=
\frac{
e^{b_{ji}+ta_{ji}}
}{
\sum_k e^{b_{ki}+ta_{ki}}
}.
\]
Then
\[
\frac{d}{dt}\pi_{j\mid i}(t)
=
\pi_{j\mid i}(t)
\left(
a_{ji}
-
\sum_k\pi_{k\mid i}(t)a_{ki}
\right).
\]
Therefore
\begin{align*}
\frac{d^2}{dt^2}\ell_i(W+tH)
&=
\sum_j
\frac{d}{dt}\pi_{j\mid i}(t)\,a_{ji}
\\
&=
\sum_j\pi_{j\mid i}(t)a_{ji}^2
-
\left(
\sum_j\pi_{j\mid i}(t)a_{ji}
\right)^2.
\end{align*}
At \(t=0\), this is exactly
\[
\Var_{j\sim\pi_{\cdot\mid i}(W)}
\left(
u_j^\top Hv_i
\right).
\]
Summing with weights \(p_i\) proves
\eqref{eq:second-directional}.
\end{proof}

\begin{proposition}[Strict convexity under two spanning conditions]
\label{prop:strict-convexity}
Assume
\begin{equation}
\Span\{u_j-u_1:2\le j\le N\}
=
\R^d
\label{eq:u-span}
\end{equation}
and
\begin{equation}
\Span\{v_i:1\le i\le N\}
=
\R^d.
\label{eq:v-span}
\end{equation}
Then:

\begin{enumerate}[label=(\roman*)]
\item the lineality space satisfies \(\cN=\{0\}\);

\item for every \(W\in\R^{d\times d}\) and every nonzero
\(H\in\R^{d\times d}\),
\[
\frac{d^2}{dt^2}L(W+tH)\bigg|_{t=0}>0;
\]

\item \(L\) is strictly convex on \(\R^{d\times d}\);

\item \(L\) has at most one finite minimizer.
\end{enumerate}
\end{proposition}

\begin{proof}
Suppose \(H\in\cN\). Then
\[
(u_j-u_1)^\top Hv_i=0
\qquad
\text{for every }i,j.
\]
Fix \(i\). By \eqref{eq:u-span}, the vectors
\[
u_j-u_1,\qquad 2\le j\le N,
\]
span \(\R^d\). Hence \(Hv_i\) is orthogonal to all of \(\R^d\), so
\[
Hv_i=0.
\]
This holds for every \(i\). Since the \(v_i\) span \(\R^d\) by
\eqref{eq:v-span}, the linear map \(H\) vanishes on a spanning set.
Therefore
\[
H=0.
\]
Thus \(\cN=\{0\}\).

Now fix \(W\) and \(H\neq 0\). By Lemma~\ref{lem:hessian},
\[
\frac{d^2}{dt^2}L(W+tH)\bigg|_{t=0}
=
\sum_i p_i
\Var_{j\sim\pi_{\cdot\mid i}(W)}
\left(
u_j^\top Hv_i
\right).
\]
Every summand is nonnegative.

Suppose the full sum were zero. Since \(p_i>0\), every variance would
have to vanish:
\[
\Var_{j\sim\pi_{\cdot\mid i}(W)}
\left(
u_j^\top Hv_i
\right)
=
0
\qquad
\text{for every }i.
\]
The distribution \(\pi_{\cdot\mid i}(W)\) has full support because all
softmax probabilities are strictly positive. A finite random variable
has zero variance if and only if it is constant on its support.
Therefore, for every fixed \(i\),
\[
u_j^\top Hv_i
=
u_1^\top Hv_i
\qquad
\text{for every }j.
\]
Equivalently,
\[
(u_j-u_1)^\top Hv_i=0
\qquad
\text{for every }i,j.
\]
Thus \(H\in\cN\). Since \(\cN=\{0\}\), this implies \(H=0\), contradicting
the assumption.

Hence the second directional derivative is strictly positive in every
nonzero direction.

To prove strict convexity, take distinct matrices \(W_0,W_1\), set
\[
H=W_1-W_0\neq 0,
\]
and define
\[
g(t)=L(W_0+tH),
\qquad
t\in[0,1].
\]
The preceding argument gives
\[
g''(t)>0
\qquad
\text{for every }t\in[0,1].
\]
Thus \(g\) is strictly convex, and consequently, for every
\(\lambda\in(0,1)\),
\[
L((1-\lambda)W_0+\lambda W_1)
<
(1-\lambda)L(W_0)+\lambda L(W_1).
\]
Therefore \(L\) is strictly convex.

A strictly convex function has at most one minimizer: if \(W_0\neq W_1\)
were both minimizers, then
\[
L\left(\frac{W_0+W_1}{2}\right)
<
\frac12L(W_0)+\frac12L(W_1),
\]
which would be strictly smaller than the minimum value. This is
impossible.
\end{proof}

\begin{lemma}[Gaussian spanning]
\label{lem:gaussian-spanning}
If \(N\ge d+1\), then the spanning conditions
\eqref{eq:u-span}--\eqref{eq:v-span} hold with probability one.
\end{lemma}

\begin{proof}
Consider first \(v_1,\dots,v_d\). The determinant
\[
\det[v_1,\dots,v_d]
\]
is a polynomial in the entries of \(v_1,\dots,v_d\). It is not the zero
polynomial, because at the deterministic choice
\[
v_i=e_i,\qquad 1\le i\le d,
\]
the determinant equals \(1\). Since the joint Gaussian law has a density
with respect to Lebesgue measure, the zero set of a nonzero polynomial
has probability zero. Therefore
\[
\det[v_1,\dots,v_d]\neq 0
\]
almost surely, and the \(v_i\) span \(\R^d\).

Similarly, consider
\[
\det[u_2-u_1,\dots,u_{d+1}-u_1].
\]
This is a polynomial in the entries of \(u_1,\dots,u_{d+1}\). It is not
identically zero: at the deterministic configuration
\[
u_1=0,
\qquad
u_{k+1}=e_k,
\qquad
1\le k\le d,
\]
the determinant equals \(1\). Hence the determinant is nonzero almost
surely, and
\[
u_2-u_1,\dots,u_{d+1}-u_1
\]
span \(\R^d\).
\end{proof}

% ============================================================
\section{Cyclic constraints and reduction to a hemisphere event}
% ============================================================

Set
\[
n=d^2.
\]
We identify \(\R^{d\times d}\) with \(\R^n\) using the vectorization map
\[
\vecop:\R^{d\times d}\to\R^{d^2}.
\]
The Frobenius inner product satisfies
\[
\ip{A}{B}
=
\vecop(A)^\top\vecop(B).
\]

Use cyclic indexing:
\[
u_{N+1}=u_1.
\]
Define
\begin{equation}
X_i
=
(u_i-u_{i+1})v_i^\top
\in\R^{d\times d},
\qquad
1\le i\le N.
\label{eq:cyclic-X}
\end{equation}

\begin{lemma}[Every separator satisfies the cyclic inequalities]
\label{lem:separator-cyclic}
If \(H\in\cS\), then
\begin{equation}
\ip{H}{X_i}\ge 0
\qquad
\text{for every }1\le i\le N.
\label{eq:cyclic-inequality}
\end{equation}
Consequently,
\begin{equation}
\{\cS\neq\{0\}\}
\subseteq
\left\{
\exists H\neq 0:
\ip{H}{X_i}\ge 0
\text{ for every }i
\right\}.
\label{eq:separator-subset}
\end{equation}
\end{lemma}

\begin{proof}
If \(H\in\cS\), then, for query \(v_i\), the correct candidate \(u_i\)
weakly dominates the particular competitor \(u_{i+1}\):
\[
u_i^\top Hv_i
\ge
u_{i+1}^\top Hv_i.
\]
Therefore
\[
(u_i-u_{i+1})^\top Hv_i\ge 0.
\]
Using the Frobenius inner product,
\begin{align*}
\ip{H}{X_i}
&=
\operatorname{Tr}
\left(
H^\top
(u_i-u_{i+1})v_i^\top
\right)
\\
&=
v_i^\top H^\top(u_i-u_{i+1})
\\
&=
(u_i-u_{i+1})^\top Hv_i.
\end{align*}
This proves \eqref{eq:cyclic-inequality}. If \(\cS\) contains a nonzero
matrix, the same matrix satisfies all cyclic inequalities, proving
\eqref{eq:separator-subset}.
\end{proof}

% ============================================================
\section{General position of the cyclic random matrices}
% ============================================================

\begin{definition}[General position]
Vectors \(x_1,\dots,x_m\in\R^n\), with \(m\ge n\), are said to be in
general position if every subset of \(n\) vectors is linearly independent.
\end{definition}

\begin{lemma}[General position of the cyclic constraints]
\label{lem:general-position}
Assume \(N>d^2\). Then the random matrices
\[
X_i=(u_i-u_{i+1})v_i^\top,
\qquad
1\le i\le N,
\]
viewed as vectors in \(\R^{d^2}\), are in general position with
probability one.
\end{lemma}

\begin{proof}
Let
\[
n=d^2.
\]
Fix a subset
\[
S\subseteq\{1,\dots,N\},
\qquad
|S|=n.
\]
Choose any ordering of \(S\) and form the \(n\times n\) matrix whose
columns are the vectorized cyclic matrices:
\[
M_S(U,V)
=
\left[
\vecop\bigl((u_i-u_{i+1})v_i^\top\bigr)
\right]_{i\in S}.
\]
Define
\[
D_S(U,V)=\det M_S(U,V).
\]
Each entry of \(M_S(U,V)\) is a polynomial in the coordinates of
\(u_1,\dots,u_N,v_1,\dots,v_N\), so \(D_S\) is a polynomial.

We prove that \(D_S\) is not identically zero by constructing one
deterministic configuration at which it is nonzero.

Because
\[
|S|=d^2,
\]
we may choose a bijection
\[
S\longleftrightarrow\{(a,b):1\le a,b\le d\}.
\]
Write \(i_{ab}\in S\) for the index corresponding to \((a,b)\).

Introduce cyclic differences
\[
\delta_i=u_i-u_{i+1},
\qquad
u_{N+1}=u_1.
\]
For each \(i_{ab}\in S\), prescribe
\[
\delta_{i_{ab}}=e_a,
\qquad
v_{i_{ab}}=e_b.
\]

Because \(N>|S|\), choose an index
\[
r\notin S.
\]
Set
\[
\delta_r
=
-\sum_{i\in S}\delta_i,
\]
and set
\[
\delta_i=0
\qquad
\text{for all }i\notin S\cup\{r\}.
\]
Then
\[
\sum_{i=1}^{N}\delta_i=0.
\]

We now verify that these prescribed cyclic differences are realizable.
Set
\[
u_1=0,
\]
and recursively define
\[
u_{i+1}=u_i-\delta_i,
\qquad
1\le i\le N-1.
\]
Then
\[
u_N
=
-\sum_{i=1}^{N-1}\delta_i.
\]
Since
\[
\sum_{i=1}^{N}\delta_i=0,
\]
we have
\[
\delta_N
=
-\sum_{i=1}^{N-1}\delta_i
=
u_N-u_1.
\]
Thus
\[
\delta_i=u_i-u_{i+1}
\]
holds for every cyclic index \(i\).

At this deterministic configuration, for \(i=i_{ab}\in S\),
\[
X_{i_{ab}}
=
\delta_{i_{ab}}v_{i_{ab}}^\top
=
e_ae_b^\top.
\]
The matrices
\[
\{e_ae_b^\top:1\le a,b\le d\}
\]
form the standard basis of \(\R^{d\times d}\). Their vectorizations form
the standard basis of \(\R^{d^2}\). Therefore
\[
D_S(U,V)=\pm 1
\]
at this configuration. In particular, \(D_S\) is not the zero polynomial.

The joint law of all Gaussian coordinates is absolutely continuous.
The zero set of a nonzero polynomial has Lebesgue measure zero, hence
\[
\Pp(D_S(U,V)=0)=0.
\]

There are only finitely many subsets \(S\) of size \(d^2\). Taking a
finite union,
\[
\Pp\left(
\exists S,\ |S|=d^2,\ D_S(U,V)=0
\right)
=0.
\]
Therefore every collection of \(d^2\) vectors among
\(\vecop(X_1),\dots,\vecop(X_N)\) is linearly independent almost surely.
\end{proof}

% ============================================================
\section{Cover's homogeneous dichotomy count}
% ============================================================

We now prove the combinatorial formula needed for the hemisphere
probability.

Let
\[
m\ge n\ge 1,
\]
and let
\[
x_1,\dots,x_m\in\R^n
\]
be in general position. A sign vector
\[
\sigma=(\sigma_1,\dots,\sigma_m)\in\{-1,+1\}^m
\]
is called homogeneously linearly realizable if there exists
\(h\in\R^n\) such that
\[
\sigma_i\ip{h}{x_i}>0
\qquad
\text{for every }i.
\]

For a fixed general-position configuration
\[
x=(x_1,\dots,x_m),
\]
let \(C_x(m,n)\) denote the number of homogeneously realizable sign
vectors.

\begin{lemma}[Cover's homogeneous counting formula]
\label{lem:cover}
Let
\[
m\ge n\ge1,
\]
and let
\[
x_1,\dots,x_m\in\R^n
\]
be in general position. Then
\begin{equation}
C_x(m,n)
=
2\sum_{k=0}^{n-1}\binom{m-1}{k}.
\label{eq:cover-formula}
\end{equation}
In particular, the number of realizable homogeneous dichotomies depends
only on \(m\) and \(n\), and not on the particular general-position
configuration. We therefore denote this common number by \(C(m,n)\).
\end{lemma}

\begin{proof}
We prove the result by strong induction on \(m+n\), over the parameter
range
\[
m\ge n\ge1.
\]

We first establish the boundary cases.

If \(n=1\), general position means that every \(x_i\) is nonzero.
A homogeneous separator \(h\in\R\setminus\{0\}\) has only two possible
orientations, corresponding to \(h>0\) and \(h<0\). Therefore
\[
C_x(m,1)=2.
\]
This agrees with
\[
2\sum_{k=0}^{0}\binom{m-1}{k}=2.
\]

Next suppose \(m=n\). General position implies that the \(n\times n\)
matrix
\[
X=[x_1,\dots,x_n]
\]
is invertible. Given any sign vector
\[
\sigma\in\{-1,+1\}^n,
\]
define
\[
h=(X^\top)^{-1}\sigma.
\]
Then
\[
X^\top h=\sigma,
\]
and hence
\[
\ip{h}{x_i}=\sigma_i
\qquad
\text{for every }i.
\]
Thus every one of the \(2^n\) sign vectors is realizable, so
\[
C_x(n,n)=2^n.
\]
Again this agrees with the claimed formula:
\[
2\sum_{k=0}^{n-1}\binom{n-1}{k}
=
2\cdot 2^{n-1}
=
2^n.
\]

It remains to consider
\[
m>n\ge2.
\]
The first \(m-1\) vectors
\[
x_1,\dots,x_{m-1}
\]
are themselves in general position in \(\R^n\).

Fix a realizable labeling
\[
(\sigma_1,\dots,\sigma_{m-1})
\]
of these first \(m-1\) vectors. The feasible set
\[
\mathcal K_\sigma
=
\left\{
h\in\R^n:
\sigma_i\ip{h}{x_i}>0
\text{ for }1\le i\le m-1
\right\}
\]
is a nonempty open subset of \(\R^n\). Since the hyperplane
\(x_m^\perp\) has empty interior, there exists
\[
h\in\mathcal K_\sigma
\]
such that
\[
\ip{h}{x_m}\neq0.
\]
The sign of \(\ip{h}{x_m}\) gives at least one realizable extension of
the labeling to \(x_m\).

Some labelings of the first \(m-1\) vectors admit both possible labels
on \(x_m\). We claim that such labelings are exactly those realizable by
a vector in \(x_m^\perp\).

Suppose first that both extensions are realizable. Then there exist
\(h_+,h_-\in\R^n\) such that
\[
\sigma_i\ip{h_+}{x_i}>0,
\qquad
\sigma_i\ip{h_-}{x_i}>0,
\qquad
1\le i\le m-1,
\]
while
\[
\ip{h_+}{x_m}>0,
\qquad
\ip{h_-}{x_m}<0.
\]
For \(t\in[0,1]\), define
\[
h_t=(1-t)h_++th_-.
\]
For every \(i\le m-1\),
\begin{align*}
\sigma_i\ip{h_t}{x_i}
&=
(1-t)\sigma_i\ip{h_+}{x_i}
+
t\sigma_i\ip{h_-}{x_i}
\\
&>0.
\end{align*}
Moreover,
\[
t\longmapsto\ip{h_t}{x_m}
\]
is continuous and changes sign. Therefore there exists
\(t_0\in(0,1)\) such that
\[
\ip{h_{t_0}}{x_m}=0.
\]
Hence the labeling is realizable by a vector in \(x_m^\perp\).

Conversely, suppose there exists
\[
h_0\in x_m^\perp
\]
such that
\[
\sigma_i\ip{h_0}{x_i}>0,
\qquad
1\le i\le m-1.
\]
Choose \(g\in\R^n\) satisfying
\[
\ip{g}{x_m}>0.
\]
Because there are only finitely many strict inequalities, there exists
\(\varepsilon_0>0\) such that, whenever
\(0<\varepsilon<\varepsilon_0\),
\[
\sigma_i\ip{h_0+\varepsilon g}{x_i}>0
\]
and
\[
\sigma_i\ip{h_0-\varepsilon g}{x_i}>0
\]
for every \(1\le i\le m-1\). Furthermore,
\[
\ip{h_0+\varepsilon g}{x_m}
=
\varepsilon\ip{g}{x_m}>0,
\]
whereas
\[
\ip{h_0-\varepsilon g}{x_m}
=
-\varepsilon\ip{g}{x_m}<0.
\]
Thus both labels on \(x_m\) are realizable.

Let
\[
P:\R^n\to x_m^\perp
\]
be the orthogonal projection. For every \(h_0\in x_m^\perp\),
\[
\ip{h_0}{x_i}
=
\ip{h_0}{Px_i}.
\]
Therefore, the labelings of \(x_1,\dots,x_{m-1}\) that admit both
extensions are exactly the homogeneously realizable labelings of
\[
Px_1,\dots,Px_{m-1}
\]
inside the \((n-1)\)-dimensional space \(x_m^\perp\).

We verify that these projected vectors are in general position in
\(x_m^\perp\). Choose any distinct indices
\[
i_1,\dots,i_{n-1}\in\{1,\dots,m-1\},
\]
and suppose
\[
\sum_{r=1}^{n-1}a_rPx_{i_r}=0.
\]
Then
\[
P\left(
\sum_{r=1}^{n-1}a_rx_{i_r}
\right)=0,
\]
so there exists a scalar \(c\) such that
\[
\sum_{r=1}^{n-1}a_rx_{i_r}=cx_m.
\]
Equivalently,
\[
\sum_{r=1}^{n-1}a_rx_{i_r}-cx_m=0.
\]
The \(n\) vectors
\[
x_{i_1},\dots,x_{i_{n-1}},x_m
\]
are linearly independent by general position. Hence
\[
a_1=\cdots=a_{n-1}=c=0.
\]
Thus every collection of \(n-1\) projected vectors is linearly
independent, proving general position in \(x_m^\perp\).

By the induction hypothesis, the number of realizable labelings of the
first \(m-1\) vectors is
\[
2\sum_{k=0}^{n-1}\binom{m-2}{k},
\]
and the number of those labelings admitting both extensions is
\[
2\sum_{k=0}^{n-2}\binom{m-2}{k}.
\]
Every realizable labeling of the first \(m-1\) vectors has one extension,
and precisely the latter labelings have one additional extension.
Therefore
\begin{align*}
C_x(m,n)
&=
2\sum_{k=0}^{n-1}\binom{m-2}{k}
+
2\sum_{k=0}^{n-2}\binom{m-2}{k}
\\
&=
2\sum_{k=0}^{n-1}
\left[
\binom{m-2}{k}
+
\binom{m-2}{k-1}
\right]
\\
&=
2\sum_{k=0}^{n-1}\binom{m-1}{k},
\end{align*}
where
\[
\binom{m-2}{-1}=0.
\]
This completes the induction.
\end{proof}

% ============================================================
\section{The Wendel--Cover hemisphere probability}
% ============================================================

\begin{lemma}[Weak and strict hemisphere events coincide]
\label{lem:weak-strict}
Let
\[
x_1,\dots,x_m\in\R^n,
\qquad
m\ge n,
\]
be in general position. Then
\[
\exists h\neq 0:
\ip{h}{x_i}\ge 0
\text{ for all }i
\]
if and only if
\[
\exists \widetilde h:
\ip{\widetilde h}{x_i}>0
\text{ for all }i.
\]
\end{lemma}

\begin{proof}
The strict event trivially implies the weak event.

For the converse, suppose \(h\neq 0\) satisfies
\[
\ip{h}{x_i}\ge 0
\qquad
\text{for all }i.
\]
Define
\[
Z=\{i:\ip{h}{x_i}=0\},
\qquad
P=\{i:\ip{h}{x_i}>0\}.
\]

Because all vectors \(x_i\) with \(i\in Z\) lie in the
\((n-1)\)-dimensional hyperplane \(h^\perp\), general position implies
\[
|Z|\le n-1.
\]
Moreover, the vectors \((x_i)_{i\in Z}\) are linearly independent.

Since they are linearly independent and \(|Z|\le n-1\), the linear
system
\[
\ip{g}{x_i}=1,
\qquad
i\in Z,
\]
has a solution \(g\in\R^n\). One way to see this is to define a linear
functional taking value \(1\) on each vector \(x_i\), extend it from
their span to all of \(\R^n\), and represent the resulting functional by
an inner product with some \(g\).

For \(i\in Z\),
\[
\ip{h+\varepsilon g}{x_i}
=
\varepsilon>0.
\]
For \(i\in P\), the quantities \(\ip{h}{x_i}\) are strictly positive.
Because \(P\) is finite, there exists \(\varepsilon_0>0\) such that,
for every \(0<\varepsilon<\varepsilon_0\),
\[
\ip{h+\varepsilon g}{x_i}>0
\qquad
\text{for every }i\in P.
\]
Thus
\[
\widetilde h=h+\varepsilon g
\]
satisfies
\[
\ip{\widetilde h}{x_i}>0
\qquad
\text{for every }i.
\]
\end{proof}

\begin{theorem}[Wendel--Cover formula]
\label{thm:wendel-cover}
Let
\[
Y_1,\dots,Y_m\in\R^n
\]
be independent random vectors. Assume:

\begin{enumerate}[label=(\roman*)]
\item each distribution is centrally symmetric:
\[
Y_i\overset{d}{=}-Y_i;
\]

\item \(Y_1,\dots,Y_m\) are in general position almost surely;

\item \(m\ge n\).
\end{enumerate}

Then
\begin{equation}
\Pp\left(
\exists h\neq 0:
\ip{h}{Y_i}\ge 0
\text{ for all }i
\right)
=
2^{1-m}
\sum_{k=0}^{n-1}\binom{m-1}{k}.
\label{eq:wendel}
\end{equation}
The random vectors need not be identically distributed.
\end{theorem}

\begin{proof}
By Lemma~\ref{lem:weak-strict}, on the almost-sure general-position event,
the weak hemisphere event is equivalent to
\[
\exists h:
\ip{h}{Y_i}>0
\qquad
\text{for every }i.
\]
It therefore suffices to compute the probability of the strict event.

Let
\[
\varepsilon_1,\dots,\varepsilon_m
\]
be independent Rademacher random variables, independent of
\(Y_1,\dots,Y_m\):
\[
\Pp(\varepsilon_i=1)
=
\Pp(\varepsilon_i=-1)
=
\frac12.
\]
By independence and central symmetry,
\[
(\varepsilon_1Y_1,\dots,\varepsilon_mY_m)
\overset{d}{=}
(Y_1,\dots,Y_m).
\]

Let \(E(y_1,\dots,y_m)\) denote the event that the deterministic vectors
\(y_1,\dots,y_m\) lie in a common open homogeneous halfspace:
\[
E(y_1,\dots,y_m)
=
\left\{
\exists h:
\ip{h}{y_i}>0
\text{ for all }i
\right\}.
\]
Then
\begin{align}
\Pp(E(Y_1,\dots,Y_m))
&=
\Pp(E(\varepsilon_1Y_1,\dots,\varepsilon_mY_m))
\nonumber\\
&=
\E_Y
\left[
2^{-m}
\sum_{\varepsilon\in\{-1,+1\}^m}
\1\{
E(\varepsilon_1Y_1,\dots,\varepsilon_mY_m)
\}
\right].
\label{eq:sign-average}
\end{align}

Fix a deterministic general-position realization
\[
y_1,\dots,y_m.
\]
For a sign vector \(\varepsilon\),
\[
E(\varepsilon_1y_1,\dots,\varepsilon_my_m)
\]
holds if and only if there exists \(h\) such that
\[
\ip{h}{\varepsilon_i y_i}>0
\qquad
\text{for all }i,
\]
or equivalently,
\[
\varepsilon_i\ip{h}{y_i}>0
\qquad
\text{for all }i.
\]
Thus the number of sign vectors \(\varepsilon\) for which the event holds
is exactly the number \(C(m,n)\) of homogeneously realizable dichotomies.

By Lemma~\ref{lem:cover},
\[
C(m,n)
=
2\sum_{k=0}^{n-1}\binom{m-1}{k}.
\]
Therefore the bracket in \eqref{eq:sign-average} equals
\[
2^{-m}C(m,n)
=
2^{1-m}
\sum_{k=0}^{n-1}\binom{m-1}{k}
\]
for every general-position realization. Since general position holds
almost surely, taking expectation yields \eqref{eq:wendel}.
\end{proof}

% ============================================================
\section{Application to the cyclic constraints}
% ============================================================

\begin{proposition}[Probability of a nonzero cyclic separator]
\label{prop:cyclic-probability}
Let
\[
X_i=(u_i-u_{i+1})v_i^\top,
\qquad
u_{N+1}=u_1.
\]
If \(N>d^2\), then
\begin{equation}
\Pp\left(
\exists H\neq 0:
\ip{H}{X_i}\ge 0
\text{ for all }i
\right)
=
2^{1-N}
\sum_{k=0}^{d^2-1}\binom{N-1}{k}.
\label{eq:cyclic-probability}
\end{equation}
\end{proposition}

\begin{proof}
Condition on
\[
U=(u_1,\dots,u_N).
\]
For fixed \(U\),
\[
X_i=(u_i-u_{i+1})v_i^\top
\]
depends only on \(v_i\). Since \(v_1,\dots,v_N\) are independent, the
conditional random matrices
\[
X_1,\dots,X_N
\]
are independent.

Moreover,
\[
v_i\overset{d}{=}-v_i,
\]
so, conditionally on \(U\),
\[
X_i
=
(u_i-u_{i+1})v_i^\top
\overset{d}{=}
-(u_i-u_{i+1})v_i^\top
=
-X_i.
\]
Thus each conditional distribution is centrally symmetric.

By Lemma~\ref{lem:general-position},
\[
\Pp_{U,V}
\left(
X_1,\dots,X_N
\text{ are not in general position}
\right)
=0.
\]
By Fubini's theorem, for almost every fixed realization \(U\),
\[
\Pp_V
\left(
X_1,\dots,X_N
\text{ are not in general position}
\mid U
\right)
=0.
\]

Therefore, for almost every \(U\), the conditional distributions satisfy
all assumptions of Theorem~\ref{thm:wendel-cover} with
\[
m=N,
\qquad
n=d^2.
\]
Hence
\[
\Pp\left(
\exists H\neq 0:
\ip{H}{X_i}\ge 0\ \forall i
\mid U
\right)
=
2^{1-N}
\sum_{k=0}^{d^2-1}\binom{N-1}{k}.
\]
The right-hand side does not depend on \(U\). Taking expectation over
\(U\) proves \eqref{eq:cyclic-probability}.
\end{proof}

\begin{corollary}[Probability of a nontrivial full separator]
\label{cor:full-separator}
If \(N>d^2\), then
\begin{equation}
\Pp(\cS\neq\{0\})
\le
2^{1-N}
\sum_{k=0}^{d^2-1}\binom{N-1}{k}.
\label{eq:full-separator-probability}
\end{equation}
\end{corollary}

\begin{proof}
By Lemma~\ref{lem:separator-cyclic},
\[
\{\cS\neq\{0\}\}
\subseteq
\left\{
\exists H\neq 0:
\ip{H}{X_i}\ge 0
\text{ for all }i
\right\}.
\]
Taking probabilities and applying
Proposition~\ref{prop:cyclic-probability} gives the result.
\end{proof}

% ============================================================
\section{Proof of the main theorem}
% ============================================================

\begin{proof}[Proof of Theorem~\ref{thm:main}]
Let
\[
\mathcal{E}_{\mathrm{span}}
=
\left\{
\Span\{u_j-u_1:2\le j\le N\}=\R^d
\right\}
\cap
\left\{
\Span\{v_i:1\le i\le N\}=\R^d
\right\}.
\]
Because \(N>d^2\ge d\), we have \(N\ge d+1\). By
Lemma~\ref{lem:gaussian-spanning},
\[
\Pp(\mathcal{E}_{\mathrm{span}})=1.
\]

Let
\[
\mathcal{E}_{\mathrm{nosep}}
=
\{\cS=\{0\}\}.
\]
By Corollary~\ref{cor:full-separator},
\[
\Pp(\mathcal{E}_{\mathrm{nosep}}^c)
=
\Pp(\cS\neq\{0\})
\le
2^{1-N}
\sum_{k=0}^{d^2-1}\binom{N-1}{k}.
\]
Therefore
\[
\Pp(\mathcal{E}_{\mathrm{nosep}})
\ge
1-
2^{1-N}
\sum_{k=0}^{d^2-1}\binom{N-1}{k}.
\]

On \(\mathcal{E}_{\mathrm{nosep}}\), Proposition
\ref{prop:deterministic-existence} shows that \(L\) is coercive and
attains its minimum at a finite matrix.

On \(\mathcal{E}_{\mathrm{span}}\), Proposition
\ref{prop:strict-convexity} shows that \(L\) is strictly convex and hence
has at most one minimizer.

Thus, on
\[
\mathcal{E}_{\mathrm{span}}
\cap
\mathcal{E}_{\mathrm{nosep}},
\]
the loss has exactly one finite minimizer. Since
\(\mathcal{E}_{\mathrm{span}}\) has probability one,
\begin{align*}
\Pp\left(
L\text{ has a unique finite minimizer}
\right)
&\ge
\Pp(\mathcal{E}_{\mathrm{nosep}})
\\
&\ge
1-
2^{1-N}
\sum_{k=0}^{d^2-1}\binom{N-1}{k}.
\end{align*}
This proves \eqref{eq:main-probability}.

It remains to derive the exponential estimate when
\[
N=\lceil d^\beta\rceil,
\qquad
\beta>2.
\]
Let
\[
B\sim\operatorname{Bin}(N-1,1/2).
\]
Then
\begin{align}
2^{1-N}
\sum_{k=0}^{d^2-1}\binom{N-1}{k}
&=
\Pp(B\le d^2-1).
\label{eq:binomial-tail-identity}
\end{align}

Because \(\beta>2\),
\[
\frac{d^2}{N}
\longrightarrow 0.
\]
Hence, for all sufficiently large \(d\),
\[
d^2-1
\le
\frac{N-1}{4}.
\]
The mean of \(B\) is
\[
\mu=\E B=\frac{N-1}{2}.
\]
Therefore
\[
d^2-1
\le
\frac{\mu}{2}.
\]
Using the multiplicative Chernoff bound
\[
\Pp(B\le(1-\delta)\mu)
\le
\exp\left(-\frac{\delta^2\mu}{2}\right),
\qquad
0<\delta<1,
\]
with \(\delta=1/2\), we obtain
\begin{align*}
\Pp(B\le d^2-1)
&\le
\Pp\left(B\le\frac{\mu}{2}\right)
\\
&\le
\exp\left(-\frac{\mu}{8}\right)
\\
&=
\exp\left(-\frac{N-1}{16}\right).
\end{align*}
Combining this with \eqref{eq:binomial-tail-identity} proves
\eqref{eq:main-exponential}.
\end{proof}

\end{document}